\begin{papercard}{egami2023}
\cf{Stages} \stagetag{L} (what may be concluded from LLM labels), \stagetag{I} (combining machine labels with a sampled human-coded subset). NeurIPS 2023 paper by computational social scientists (Princeton, Cornell). Added from the gap search of 23 October 2026 because the thesis reports counts that partly rest on model output.

\cf{Problem} Researchers increasingly count or regress on labels produced by an LLM (``surrogates''). Using them directly gives biased estimates and wrong confidence intervals even when the labels look accurate \pc{egami2023}{1}{We show that direct use of surrogate labels in downstream statistical analyses leads to substantial bias and invalid confidence intervals, even with high surrogate accuracy of 80–90\%}. The reason is that errors are not random with respect to the variables of interest.

\cf{Method} Design-based supervised learning (DSL): an expert labels a small sample drawn with known probabilities, and a doubly robust estimator combines those gold labels with the surrogates \pc{egami2023}{1}{DSL employs a doubly-robust procedure to combine surrogate labels with a smaller number of high-quality, gold-standard labels}. Validity comes from the sampling design, not from the quality of the model \pc{egami2023}{1}{even when surrogates are arbitrarily biased and without requiring stringent assumptions, by controlling the probability of sampling documents for gold-standard labeling}. The only requirements are known sampling probabilities bounded away from zero \pc{egami2023}{3}{These strong theoretical guarantees only require that the sampling probability of the documents selected for gold-standard be known and be bounded away from zero}; better surrogates make the estimate more precise, not more valid.

\cf{Limits (card)} Aimed at regression and proportion estimates on large corpora; the thesis is descriptive with 50 apps, and a design-based correction may be more than it needs. The gold labels are assumed correct, which a single coder cannot guarantee.

\cf{Use for this thesis} (1) The principle behind the methodology's rule that model labels are screened, not counted: any prevalence statement (``$x$ of 50 apps name harm $Y$'') must rest on human-checked labels, or on a correction like DSL. (2) Positives are all reviewed in the current design; the open risk is model negatives, i.e.\ absence claims. The missed-mention audit samples units outside the vocabulary hits; if its sampling probabilities are recorded, the same logic yields an estimate of missed mentions instead of a bare count. (3) Write the sampling probabilities of the validation and audit samples into the protocol before drawing them.
\end{papercard}
